\section{Introduction}
Whole-slide images (WSIs) capture tissue morphology over large spatial extents at high resolution. They support tumor measurement, tissue-composition analysis, and region-of-interest selection. Weakly supervised tissue segmenters such as HistoSegNet and CAMEL reduce annotation requirements while predicting predefined tissue categories~\cite{chan2019histosegnet,xu2019camel}. Complex tissue boundaries and substantial scale variation make dense pixel annotation expensive. The intended target also changes with the study: the same slide may require tumor epithelium, stroma, or adipose tissue to be extracted separately. A fixed set of segmentation channels cannot directly express this user-defined target.

Click-based interactive methods and promptable foundation models provide a direct way to specify what to segment~\cite{sofiiuk2022ritm,liu2023simpleclick,kirillov2023sam}. Applying it to an entire WSI, however, presents three challenges. First, the same tissue may occupy disconnected regions, requiring a local prompt to convey information beyond its neighborhood. Second, different tissues can share local appearance, so target identification depends on surrounding architecture and examples of what to exclude. Third, recognizing tissue requires broad context, whereas delineating small structures requires local detail~\cite{tokunaga2019awmf,vanrijthoven2021hooknet,ho2021dmmn}. At slide scale, these challenges are compounded by repeated prompting and image computation.

Few-click WSI segmentation must recover a target across spatially separated regions while limiting false detections where it is absent. This coverage--false-positive tradeoff becomes central when one support example guides predictions across an entire slide. FewClick combines signed support evidence with dense probability transfer to address these two demands. Whole-patch and tiled features retain tissue architecture and local detail, while multiscale region context supports target identification and correspondence links support evidence to remote query locations.

Figure~\ref{fig:overview} summarizes FewClick. A shared backbone encodes aligned fine, intermediate, and coarse views, combining whole-patch and tiled fine features. Region pooling and cross-scale messages organize tissue context. Signed support clicks condition the initial prediction; coarse-to-fine correspondence then transfers predicted support probabilities to refine query logits. This support example is reused across the slide without further human prompts.

Our contributions are threefold. (1) We introduce a whole-slide tissue propagation framework that uses a few signed clicks in one support patch, with no further human prompts in query regions. (2) We combine multiscale whole-patch/tiled features, signed support evidence, and coarse-to-fine probability transfer to address target coverage and non-target rejection, using region context to organize spatial information. (3) We evaluate six methods on 15 WSIs and 11,000 remote development tasks under common support prompts and additional adaptation budgets, and quantify the coverage--false-positive tradeoff through matched controls. FewClick achieves the highest full-slide macro Dice among the compared methods (43.75\%), exceeding SANSA by 10.52 points. On the remote development queue, it achieves the lowest target-absent FPA (7.75\%) while retaining the highest target-present macro Dice (42.11\%).
